\section{总结与延伸}

\subsection{讲者的核心总结}

\begin{figure}[H]
\centering
\includegraphics[width=0.95\textwidth]{slides-images/slide-050.jpg}
\caption{全课总结\protect\footnotemark}
\end{figure}
\footnotetext{Slides 第51页（最后一页）。}

Tatsu Hashimoto 在课程结尾强调了三点：
\begin{enumerate}
    \item \textbf{硬件驱动规模化}：底层硬件细节决定了什么能规模化、什么不能
    \item \textbf{以 matmul + 数据搬运为中心思考}：当前基于 GPU 的计算体系，要求算法设计者时刻关注矩阵乘法和数据移动
    \item \textbf{细致的 GPU 优化带来真正的性能提升}：Coalescing、Tiling、Fusion 等看似底层的技巧，是 FlashAttention 等突破性算法的基石
\end{enumerate}

\subsection{全课知识图谱}

本课建立了一条从硬件到算法的完整认知链：

\begin{center}
\begin{tikzpicture}[node distance=0.8cm, every node/.style={draw, rounded corners, minimum width=3cm, minimum height=0.7cm, font=\small}]
\node (hw) {GPU 硬件架构};
\node (mem) [below=of hw] {内存层次与瓶颈};
\node (opt) [below=of mem] {六大优化技巧};
\node (flash) [below=of opt] {FlashAttention};
\draw[->, thick] (hw) -- (mem);
\draw[->, thick] (mem) -- (opt);
\draw[->, thick] (opt) -- (flash);
\node[draw=none, right=1cm of hw, text width=5cm, font=\scriptsize] {SM/SP、Thread/Block/Warp、SIMT};
\node[draw=none, right=1cm of mem, text width=5cm, font=\scriptsize] {SRAM vs HBM、Roofline Model、Memory Wall};
\node[draw=none, right=1cm of opt, text width=5cm, font=\scriptsize] {Divergence、Low Precision、Fusion、Recompute、Coalescing、Tiling};
\node[draw=none, right=1cm of flash, text width=5cm, font=\scriptsize] {Tiled KQV + Online Softmax + Recomputation};
\end{tikzpicture}
\end{center}

\subsection{关键 Takeaways}

\begin{importantbox}{五条核心原则}
\begin{enumerate}
    \item \textbf{内存是新的瓶颈}：现代 GPU 上，计算速度远快于内存带宽，大多数优化围绕减少内存访问展开
    \item \textbf{Matmul 是``特权运算''}：尽量让你的工作负载以 matmul 为主，避免大量非 matmul 操作
    \item \textbf{矩阵维度要``友好''}：选择 64/128/256 的倍数，避免质数维度
    \item \textbf{数据要尽量留在 SRAM}：通过 Tiling 将工作数据加载到 Shared Memory，最小化全局内存访问
    \item \textbf{FlashAttention 不是魔法}：它是 Tiling + Fusion + Online Softmax + Recomputation 的精妙组合
\end{enumerate}
\end{importantbox}

\end{document}
